% Describing radiation shielding effect with lead

A un metro la Fuente nos dará: $53.5 \times 0.55 = 29.425 \, \text{R/h}$

Sabemos que $0.2"$ de Plomo reduce a la mitad el nivel de Radiación:

\begin{table}[h]
\centering
\begin{tabular}{|c|c|}
\hline
\textbf{ESPESOR} & \textbf{NIVEL RAD.} \\
\hline
0 & 29425 \, \text{mR/h} \\
0.2" & 14712.5 \, \text{mR/h} \\
0.4" & 7356.25 \, \text{mR/h} \\
0.6" & 3678.13 \, \text{mR/h} \\
0.8" & 1839.1 \, \text{mR/h} \\
1.0" & 919.5 \, \text{mR/h} \\
1.2" & 459.77 \, \text{mR/h} \\
1.4" & 229.88 \, \text{mR/h} \\
1.6" & 114.94 \, \text{mR/h} \\
1.8" & 57.47 \, \text{mR/h} \\
\hline
\end{tabular}
\caption{Reducción de la radiación con diferentes espesores de plomo}
\end{table}

% Applying inverse square law for radiation level calculation

¿Qué nivel de Radiación nos dará a 3 metros una fuente de Ir-192 de 85 Ci?

A un metro tendríamos: $0.55 \times 85 = 46.75 \, \text{R/h}$

Aplicando la Ley de Cuadrado Inverso:

\textbf{DATOS:}
\[
I_1 = 46.75 \, \text{R/h}
\]
\[
d_1 = 1 \, \text{m}
\]
\[
d_2 = 3 \, \text{m}
\]

\textbf{FÓRMULA:}
\[
I_2 = \left(\frac{d_1}{d_2}\right)^2 \times I_1 = \left(\frac{1}{3}\right)^2 \times 46.75 = 5.19 \, \text{R/h}
\]

% Scenario of a worker handling a radioactive source

A un trabajador se le cae una Fuente de 60 Ci y decide tomarla del contenedor con la mano y meterla al contenedor. Si el contenedor está a 10 cm. de la Fuente, ¿cuál será el nivel de Radiación en ese punto?

% Calculation based on provided data
- A un metro tendríamos: $0.55 \times 60 = 33 \, \text{R/h}$

% End of the radiation safety manual text portion